\section{微扰论和线性响应}

本节按 Giuliani--Vignale 的路线，从含时微扰推导 Kubo 公式。这是密度响应、介电函数与 Lindhard 函数的共同基础。

\subsection{含时微扰}

无扰哈密顿量为 $H_0$，外场通过算符 $\hat A$ 耦合：
\begin{equation}
  H(t)=H_0+H'(t),
  \qquad
  H'(t)=-\hat A\,F(t).
\end{equation}
更一般地，可写
\begin{equation}
  H'(t)
  =-\int\mathrm{d}\mathbf{r}\,
  \hat A(\mathbf{r})\,F(\mathbf{r},t).
\end{equation}
对密度响应，取 $\hat A(\mathbf{r})=\hat n(\mathbf{r})$，$F=-e\phi_{\mathrm{ext}}$（或直接把扰动势并入约定）。

薛定谔绘景中态满足 $i\hbar\partial_t|\psi\rangle=H(t)|\psi\rangle$。令
\begin{equation}
  |\psi(t)\rangle
  =e^{-\mathrm{i}H_0 t/\hbar}|\tilde\psi(t)\rangle,
\end{equation}
则相互作用绘景中
\begin{equation}
  i\hbar\partial_t|\tilde\psi\rangle
  =H'_I(t)|\tilde\psi\rangle,
  \qquad
  H'_I(t)=e^{\mathrm{i}H_0 t/\hbar}H'(t)e^{-\mathrm{i}H_0 t/\hbar}.
\end{equation}

\subsection{一阶演化与 Kubo 公式}

设 $t\to-\infty$ 时体系处于 $H_0$ 的热平衡（或基态），密度矩阵为 $\rho_0$。一阶
\begin{equation}
  |\tilde\psi(t)\rangle
  \approx|\tilde\psi(-\infty)\rangle
  -\frac{\mathrm{i}}{\hbar}
  \int_{-\infty}^{t}\mathrm{d}t'\,
  H'_I(t')|\tilde\psi(-\infty)\rangle.
\end{equation}
可观测量 $\hat B$ 的期望值相对无扰值的一阶改变为
\begin{equation}
  \delta\langle\hat B(t)\rangle
  =-\frac{\mathrm{i}}{\hbar}
  \int_{-\infty}^{t}\mathrm{d}t'\,
  \bigl\langle
  \bigl[\hat B_I(t),\hat A_I(t')\bigr]
  \bigr\rangle_0
  F(t'),
\end{equation}
其中 $\langle\cdot\rangle_0=\mathrm{Tr}(\rho_0\,\cdot)$。定义 \textbf{retard 响应函数}
\begin{equation}
  \chi_{BA}(t-t')
  =-\frac{\mathrm{i}}{\hbar}
  \theta(t-t')
  \bigl\langle
  \bigl[\hat B_I(t),\hat A_I(t')\bigr]
  \bigr\rangle_0,
\label{eq:kubo}
\end{equation}
则
\begin{equation}
  \delta\langle B(t)\rangle
  =\int_{-\infty}^{\infty}\mathrm{d}t'\,
  \chi_{BA}(t-t')F(t').
\end{equation}
式 \eqref{eq:kubo} 即 Kubo 公式：线性响应完全由平衡态推迟对易子决定。

\subsection{频率空间}

Fourier 变换
\begin{equation}
  \chi_{BA}(\omega)
  =\int_0^{\infty}\mathrm{d}t\,
  e^{\mathrm{i}\omega t}
  \Bigl(
  -\frac{\mathrm{i}}{\hbar}
  \bigr)
  \bigl\langle[\hat B(t),\hat A(0)]\bigr\rangle_0,
\end{equation}
其中已令 $\omega\to\omega+\mathrm{i}0^+$ 以保证收敛。空间非均匀时推广为 $\chi_{BA}(\mathbf{r},\mathbf{r}';\omega)$，平移不变体系下为 $\chi_{BA}(\mathbf{q},\omega)$。

\begin{note}
Kubo 公式对任意多体 $H_0$ 成立，不要求无相互作用。近似出现在如何计算平衡对易子关联——例如用自由粒子（Lindhard）或 RPA 屏蔽等。
\end{note}
